% Options for packages loaded elsewhere
\PassOptionsToPackage{unicode}{hyperref}
\PassOptionsToPackage{hyphens}{url}
\documentclass[
]{article}
\usepackage{xcolor}
\usepackage[margin=1in]{geometry}
\usepackage{amsmath,amssymb}
\setcounter{secnumdepth}{-\maxdimen} % remove section numbering
\usepackage{iftex}
\ifPDFTeX
  \usepackage[T1]{fontenc}
  \usepackage[utf8]{inputenc}
  \usepackage{textcomp} % provide euro and other symbols
\else % if luatex or xetex
  \usepackage{unicode-math} % this also loads fontspec
  \defaultfontfeatures{Scale=MatchLowercase}
  \defaultfontfeatures[\rmfamily]{Ligatures=TeX,Scale=1}
\fi
\usepackage{lmodern}
\ifPDFTeX\else
  % xetex/luatex font selection
\fi
% Use upquote if available, for straight quotes in verbatim environments
\IfFileExists{upquote.sty}{\usepackage{upquote}}{}
\IfFileExists{microtype.sty}{% use microtype if available
  \usepackage[]{microtype}
  \UseMicrotypeSet[protrusion]{basicmath} % disable protrusion for tt fonts
}{}
\makeatletter
\@ifundefined{KOMAClassName}{% if non-KOMA class
  \IfFileExists{parskip.sty}{%
    \usepackage{parskip}
  }{% else
    \setlength{\parindent}{0pt}
    \setlength{\parskip}{6pt plus 2pt minus 1pt}}
}{% if KOMA class
  \KOMAoptions{parskip=half}}
\makeatother
\usepackage{color}
\usepackage{fancyvrb}
\newcommand{\VerbBar}{|}
\newcommand{\VERB}{\Verb[commandchars=\\\{\}]}
\DefineVerbatimEnvironment{Highlighting}{Verbatim}{commandchars=\\\{\}}
% Add ',fontsize=\small' for more characters per line
\usepackage{framed}
\definecolor{shadecolor}{RGB}{248,248,248}
\newenvironment{Shaded}{\begin{snugshade}}{\end{snugshade}}
\newcommand{\AlertTok}[1]{\textcolor[rgb]{0.94,0.16,0.16}{#1}}
\newcommand{\AnnotationTok}[1]{\textcolor[rgb]{0.56,0.35,0.01}{\textbf{\textit{#1}}}}
\newcommand{\AttributeTok}[1]{\textcolor[rgb]{0.13,0.29,0.53}{#1}}
\newcommand{\BaseNTok}[1]{\textcolor[rgb]{0.00,0.00,0.81}{#1}}
\newcommand{\BuiltInTok}[1]{#1}
\newcommand{\CharTok}[1]{\textcolor[rgb]{0.31,0.60,0.02}{#1}}
\newcommand{\CommentTok}[1]{\textcolor[rgb]{0.56,0.35,0.01}{\textit{#1}}}
\newcommand{\CommentVarTok}[1]{\textcolor[rgb]{0.56,0.35,0.01}{\textbf{\textit{#1}}}}
\newcommand{\ConstantTok}[1]{\textcolor[rgb]{0.56,0.35,0.01}{#1}}
\newcommand{\ControlFlowTok}[1]{\textcolor[rgb]{0.13,0.29,0.53}{\textbf{#1}}}
\newcommand{\DataTypeTok}[1]{\textcolor[rgb]{0.13,0.29,0.53}{#1}}
\newcommand{\DecValTok}[1]{\textcolor[rgb]{0.00,0.00,0.81}{#1}}
\newcommand{\DocumentationTok}[1]{\textcolor[rgb]{0.56,0.35,0.01}{\textbf{\textit{#1}}}}
\newcommand{\ErrorTok}[1]{\textcolor[rgb]{0.64,0.00,0.00}{\textbf{#1}}}
\newcommand{\ExtensionTok}[1]{#1}
\newcommand{\FloatTok}[1]{\textcolor[rgb]{0.00,0.00,0.81}{#1}}
\newcommand{\FunctionTok}[1]{\textcolor[rgb]{0.13,0.29,0.53}{\textbf{#1}}}
\newcommand{\ImportTok}[1]{#1}
\newcommand{\InformationTok}[1]{\textcolor[rgb]{0.56,0.35,0.01}{\textbf{\textit{#1}}}}
\newcommand{\KeywordTok}[1]{\textcolor[rgb]{0.13,0.29,0.53}{\textbf{#1}}}
\newcommand{\NormalTok}[1]{#1}
\newcommand{\OperatorTok}[1]{\textcolor[rgb]{0.81,0.36,0.00}{\textbf{#1}}}
\newcommand{\OtherTok}[1]{\textcolor[rgb]{0.56,0.35,0.01}{#1}}
\newcommand{\PreprocessorTok}[1]{\textcolor[rgb]{0.56,0.35,0.01}{\textit{#1}}}
\newcommand{\RegionMarkerTok}[1]{#1}
\newcommand{\SpecialCharTok}[1]{\textcolor[rgb]{0.81,0.36,0.00}{\textbf{#1}}}
\newcommand{\SpecialStringTok}[1]{\textcolor[rgb]{0.31,0.60,0.02}{#1}}
\newcommand{\StringTok}[1]{\textcolor[rgb]{0.31,0.60,0.02}{#1}}
\newcommand{\VariableTok}[1]{\textcolor[rgb]{0.00,0.00,0.00}{#1}}
\newcommand{\VerbatimStringTok}[1]{\textcolor[rgb]{0.31,0.60,0.02}{#1}}
\newcommand{\WarningTok}[1]{\textcolor[rgb]{0.56,0.35,0.01}{\textbf{\textit{#1}}}}
\usepackage{longtable,booktabs,array}
\usepackage{calc} % for calculating minipage widths
% Correct order of tables after \paragraph or \subparagraph
\usepackage{etoolbox}
\makeatletter
\patchcmd\longtable{\par}{\if@noskipsec\mbox{}\fi\par}{}{}
\makeatother
% Allow footnotes in longtable head/foot
\IfFileExists{footnotehyper.sty}{\usepackage{footnotehyper}}{\usepackage{footnote}}
\makesavenoteenv{longtable}
\usepackage{graphicx}
\makeatletter
\newsavebox\pandoc@box
\newcommand*\pandocbounded[1]{% scales image to fit in text height/width
  \sbox\pandoc@box{#1}%
  \Gscale@div\@tempa{\textheight}{\dimexpr\ht\pandoc@box+\dp\pandoc@box\relax}%
  \Gscale@div\@tempb{\linewidth}{\wd\pandoc@box}%
  \ifdim\@tempb\p@<\@tempa\p@\let\@tempa\@tempb\fi% select the smaller of both
  \ifdim\@tempa\p@<\p@\scalebox{\@tempa}{\usebox\pandoc@box}%
  \else\usebox{\pandoc@box}%
  \fi%
}
% Set default figure placement to htbp
\def\fps@figure{htbp}
\makeatother
\setlength{\emergencystretch}{3em} % prevent overfull lines
\providecommand{\tightlist}{%
  \setlength{\itemsep}{0pt}\setlength{\parskip}{0pt}}
\usepackage{bookmark}
\IfFileExists{xurl.sty}{\usepackage{xurl}}{} % add URL line breaks if available
\urlstyle{same}
\hypersetup{
  pdftitle={DGP Scenario Reproduction},
  pdfauthor={Arjun Sekhar},
  hidelinks,
  pdfcreator={LaTeX via pandoc}}

\title{DGP Scenario Reproduction}
\author{Arjun Sekhar}
\date{2026-06-14}

\begin{document}
\maketitle

\subsection{Chunk 1: Libraries}\label{chunk-1-libraries}

\begin{Shaded}
\begin{Highlighting}[]
\FunctionTok{library}\NormalTok{(randomForestSRC)}
\FunctionTok{library}\NormalTok{(survival)}
\FunctionTok{library}\NormalTok{(PRROC)}
\FunctionTok{library}\NormalTok{(pander)}
\end{Highlighting}
\end{Shaded}

\subsection{Chunk 2: Function
definitions}\label{chunk-2-function-definitions}

\begin{Shaded}
\begin{Highlighting}[]
\FunctionTok{set.seed}\NormalTok{(}\DecValTok{44940076}\NormalTok{)}

\DocumentationTok{\#\# Step 1: Data generating process}
\NormalTok{generate\_survival\_data }\OtherTok{\textless{}{-}} \ControlFlowTok{function}\NormalTok{(n, }\AttributeTok{p =} \DecValTok{10}\NormalTok{, betas,}
                                   \AttributeTok{interactions =} \ConstantTok{NULL}\NormalTok{,}
                                   \AttributeTok{censoring\_rate =} \FloatTok{0.99}\NormalTok{) \{}
  
  \DocumentationTok{\#\# Step 1a: Generate orthogonal covariates (mimics PCA structure of ULB V1{-}V28)}
\NormalTok{  X }\OtherTok{\textless{}{-}} \FunctionTok{matrix}\NormalTok{(}\FunctionTok{rnorm}\NormalTok{(n }\SpecialCharTok{*}\NormalTok{ p), }\AttributeTok{nrow =}\NormalTok{ n, }\AttributeTok{ncol =}\NormalTok{ p)}
  \FunctionTok{colnames}\NormalTok{(X) }\OtherTok{\textless{}{-}} \FunctionTok{paste0}\NormalTok{(}\StringTok{"X"}\NormalTok{, }\DecValTok{1}\SpecialCharTok{:}\NormalTok{p)}
  
  \DocumentationTok{\#\# Step 1b: Compute linear predictor}
\NormalTok{  lp }\OtherTok{\textless{}{-}}\NormalTok{ X }\SpecialCharTok{\%*\%}\NormalTok{ betas}
  
  \DocumentationTok{\#\# Step 1c: Add interaction terms if specified (Scenario B only)}
  \ControlFlowTok{if}\NormalTok{ (}\SpecialCharTok{!}\FunctionTok{is.null}\NormalTok{(interactions)) \{}
    \ControlFlowTok{for}\NormalTok{ (term }\ControlFlowTok{in}\NormalTok{ interactions) \{}
\NormalTok{      lp }\OtherTok{\textless{}{-}}\NormalTok{ lp }\SpecialCharTok{+}\NormalTok{ term}\SpecialCharTok{$}\NormalTok{gamma }\SpecialCharTok{*}\NormalTok{ X[, term}\SpecialCharTok{$}\NormalTok{i] }\SpecialCharTok{*}\NormalTok{ X[, term}\SpecialCharTok{$}\NormalTok{j]}
\NormalTok{    \}}
\NormalTok{  \}}
  
  \DocumentationTok{\#\# Step 1d: Generate exponential survival times}
\NormalTok{  lambda     }\OtherTok{\textless{}{-}} \FunctionTok{exp}\NormalTok{(lp)}
\NormalTok{  time\_event }\OtherTok{\textless{}{-}} \FunctionTok{rexp}\NormalTok{(n, }\AttributeTok{rate =}\NormalTok{ lambda)}
  
  \DocumentationTok{\#\# Step 1e: Apply censoring to hit target censoring rate}
\NormalTok{  lambda\_c   }\OtherTok{\textless{}{-}} \FunctionTok{quantile}\NormalTok{(time\_event, }\AttributeTok{probs =} \DecValTok{1} \SpecialCharTok{{-}}\NormalTok{ censoring\_rate)}
\NormalTok{  time\_censor }\OtherTok{\textless{}{-}} \FunctionTok{rexp}\NormalTok{(n, }\AttributeTok{rate =} \DecValTok{1} \SpecialCharTok{/} \FunctionTok{as.numeric}\NormalTok{(lambda\_c))}
  
  \DocumentationTok{\#\# Step 1f: Observed time and event indicator}
\NormalTok{  time\_obs }\OtherTok{\textless{}{-}} \FunctionTok{pmin}\NormalTok{(time\_event, time\_censor)}
\NormalTok{  status   }\OtherTok{\textless{}{-}} \FunctionTok{as.integer}\NormalTok{(time\_event }\SpecialCharTok{\textless{}=}\NormalTok{ time\_censor)}
  
  \FunctionTok{data.frame}\NormalTok{(X, }\AttributeTok{time =}\NormalTok{ time\_obs, }\AttributeTok{status =}\NormalTok{ status)}
\NormalTok{\}}

\DocumentationTok{\#\# Step 2: Scenario runner}
\NormalTok{run\_scenario }\OtherTok{\textless{}{-}} \ControlFlowTok{function}\NormalTok{(scenario\_label, n, betas,}
                         \AttributeTok{interactions =} \ConstantTok{NULL}\NormalTok{,}
\NormalTok{                         censoring\_rate,}
                         \AttributeTok{n\_splits =} \DecValTok{30}\NormalTok{) \{}
  
\NormalTok{  auprc\_lr  }\OtherTok{\textless{}{-}} \FunctionTok{numeric}\NormalTok{(n\_splits)}
\NormalTok{  auprc\_rsf }\OtherTok{\textless{}{-}} \FunctionTok{numeric}\NormalTok{(n\_splits)}
\NormalTok{  cerr\_lr   }\OtherTok{\textless{}{-}} \FunctionTok{numeric}\NormalTok{(n\_splits)}
\NormalTok{  cerr\_rsf  }\OtherTok{\textless{}{-}} \FunctionTok{numeric}\NormalTok{(n\_splits)}
\NormalTok{  n\_events  }\OtherTok{\textless{}{-}} \FunctionTok{numeric}\NormalTok{(n\_splits)}
  
  \ControlFlowTok{for}\NormalTok{ (i }\ControlFlowTok{in} \FunctionTok{seq\_len}\NormalTok{(n\_splits)) \{}
    \FunctionTok{message}\NormalTok{(}\StringTok{"  ["}\NormalTok{, scenario\_label, }\StringTok{"] Split "}\NormalTok{, i, }\StringTok{" of "}\NormalTok{, n\_splits,}
            \StringTok{" — "}\NormalTok{, }\FunctionTok{Sys.time}\NormalTok{())}
    
    \DocumentationTok{\#\# Step 2a: Generate data from known DGP}
\NormalTok{    dat }\OtherTok{\textless{}{-}} \FunctionTok{generate\_survival\_data}\NormalTok{(n, }\AttributeTok{betas =}\NormalTok{ betas,}
                                  \AttributeTok{interactions =}\NormalTok{ interactions,}
                                  \AttributeTok{censoring\_rate =}\NormalTok{ censoring\_rate)}
    
\NormalTok{    n\_events[i] }\OtherTok{\textless{}{-}} \FunctionTok{sum}\NormalTok{(dat}\SpecialCharTok{$}\NormalTok{status)}
    
    \DocumentationTok{\#\# Step 2b: Stratified 70/30 train{-}test split}
\NormalTok{    event\_idx    }\OtherTok{\textless{}{-}} \FunctionTok{which}\NormalTok{(dat}\SpecialCharTok{$}\NormalTok{status }\SpecialCharTok{==} \DecValTok{1}\NormalTok{)}
\NormalTok{    nonevent\_idx }\OtherTok{\textless{}{-}} \FunctionTok{which}\NormalTok{(dat}\SpecialCharTok{$}\NormalTok{status }\SpecialCharTok{==} \DecValTok{0}\NormalTok{)}
    
\NormalTok{    train\_idx }\OtherTok{\textless{}{-}} \FunctionTok{c}\NormalTok{(}
      \FunctionTok{sample}\NormalTok{(event\_idx,    }\AttributeTok{size =} \FunctionTok{floor}\NormalTok{(}\FloatTok{0.7} \SpecialCharTok{*} \FunctionTok{length}\NormalTok{(event\_idx))),}
      \FunctionTok{sample}\NormalTok{(nonevent\_idx, }\AttributeTok{size =} \FunctionTok{floor}\NormalTok{(}\FloatTok{0.7} \SpecialCharTok{*} \FunctionTok{length}\NormalTok{(nonevent\_idx)))}
\NormalTok{    )}
    
\NormalTok{    train }\OtherTok{\textless{}{-}}\NormalTok{ dat[train\_idx, ]}
\NormalTok{    test  }\OtherTok{\textless{}{-}}\NormalTok{ dat[}\SpecialCharTok{{-}}\NormalTok{train\_idx, ]}
    
    \DocumentationTok{\#\# Step 2c: Logistic Regression}
\NormalTok{    lr\_mod  }\OtherTok{\textless{}{-}} \FunctionTok{glm}\NormalTok{(status }\SpecialCharTok{\textasciitilde{}}\NormalTok{ . }\SpecialCharTok{{-}}\NormalTok{ time, }\AttributeTok{data =}\NormalTok{ train, }\AttributeTok{family =}\NormalTok{ binomial)}
\NormalTok{    lr\_prob }\OtherTok{\textless{}{-}} \FunctionTok{predict}\NormalTok{(lr\_mod, }\AttributeTok{newdata =}\NormalTok{ test, }\AttributeTok{type =} \StringTok{"response"}\NormalTok{)}
    
    \DocumentationTok{\#\# Step 2d: Random Survival Forest}
\NormalTok{    rsf\_mod  }\OtherTok{\textless{}{-}} \FunctionTok{rfsrc}\NormalTok{(}\FunctionTok{Surv}\NormalTok{(time, status) }\SpecialCharTok{\textasciitilde{}}\NormalTok{ ., }\AttributeTok{data =}\NormalTok{ train,}
                      \AttributeTok{ntree =} \DecValTok{100}\NormalTok{, }\AttributeTok{nodesize =} \DecValTok{15}\NormalTok{,}
                      \AttributeTok{splitrule =} \StringTok{"logrank"}\NormalTok{, }\AttributeTok{fast =} \ConstantTok{TRUE}\NormalTok{, }\AttributeTok{nthread =} \DecValTok{4}\NormalTok{)}
    
\NormalTok{    rsf\_pred }\OtherTok{\textless{}{-}} \FunctionTok{predict}\NormalTok{(rsf\_mod, }\AttributeTok{newdata =}\NormalTok{ test[, rsf\_mod}\SpecialCharTok{$}\NormalTok{xvar.names,}
                                                \AttributeTok{drop =} \ConstantTok{FALSE}\NormalTok{])}
\NormalTok{    rsf\_risk }\OtherTok{\textless{}{-}} \DecValTok{1} \SpecialCharTok{{-}}\NormalTok{ rsf\_pred}\SpecialCharTok{$}\NormalTok{survival[, }\FunctionTok{ncol}\NormalTok{(rsf\_pred}\SpecialCharTok{$}\NormalTok{survival)]}
\NormalTok{    rsf\_risk }\OtherTok{\textless{}{-}} \FunctionTok{as.numeric}\NormalTok{(rsf\_risk)}
    
\NormalTok{    true\_labels }\OtherTok{\textless{}{-}}\NormalTok{ test}\SpecialCharTok{$}\NormalTok{status}
    
    \DocumentationTok{\#\# Step 2e: AUPRC}
\NormalTok{    lr\_keep  }\OtherTok{\textless{}{-}} \FunctionTok{is.finite}\NormalTok{(lr\_prob)  }\SpecialCharTok{\&} \FunctionTok{is.finite}\NormalTok{(true\_labels)}
\NormalTok{    rsf\_keep }\OtherTok{\textless{}{-}} \FunctionTok{is.finite}\NormalTok{(rsf\_risk) }\SpecialCharTok{\&} \FunctionTok{is.finite}\NormalTok{(true\_labels)}
    
\NormalTok{    auprc\_lr[i] }\OtherTok{\textless{}{-}} \FunctionTok{pr.curve}\NormalTok{(}
      \AttributeTok{scores.class0 =}\NormalTok{ lr\_prob[lr\_keep][true\_labels[lr\_keep] }\SpecialCharTok{==} \DecValTok{1}\NormalTok{],}
      \AttributeTok{scores.class1 =}\NormalTok{ lr\_prob[lr\_keep][true\_labels[lr\_keep] }\SpecialCharTok{==} \DecValTok{0}\NormalTok{],}
      \AttributeTok{curve =} \ConstantTok{FALSE}\NormalTok{)}\SpecialCharTok{$}\NormalTok{auc.integral}
    
\NormalTok{    auprc\_rsf[i] }\OtherTok{\textless{}{-}} \FunctionTok{pr.curve}\NormalTok{(}
      \AttributeTok{scores.class0 =}\NormalTok{ rsf\_risk[rsf\_keep][true\_labels[rsf\_keep] }\SpecialCharTok{==} \DecValTok{1}\NormalTok{],}
      \AttributeTok{scores.class1 =}\NormalTok{ rsf\_risk[rsf\_keep][true\_labels[rsf\_keep] }\SpecialCharTok{==} \DecValTok{0}\NormalTok{],}
      \AttributeTok{curve =} \ConstantTok{FALSE}\NormalTok{)}\SpecialCharTok{$}\NormalTok{auc.integral}
    
    \DocumentationTok{\#\# Step 2f: Concordance error (1 {-} C{-}index)}
    \DocumentationTok{\#\# Note: concordance() uses the survival convention where higher predictor =}
    \DocumentationTok{\#\# longer survival. Since lr\_prob and rsf\_risk are RISK scores (higher = sooner}
    \DocumentationTok{\#\# failure), the raw concordance returns \textless{} 0.5. CErr = 1 {-} concordance()}
    \DocumentationTok{\#\# therefore recovers the true discriminating C{-}index (\textgreater{} 0.5 = good).}
    \DocumentationTok{\#\# This must be stated explicitly when reporting these values in the thesis.}
\NormalTok{    cerr\_lr[i]  }\OtherTok{\textless{}{-}} \DecValTok{1} \SpecialCharTok{{-}} \FunctionTok{as.numeric}\NormalTok{(}\FunctionTok{concordance}\NormalTok{(}
      \FunctionTok{Surv}\NormalTok{(test}\SpecialCharTok{$}\NormalTok{time, test}\SpecialCharTok{$}\NormalTok{status) }\SpecialCharTok{\textasciitilde{}}\NormalTok{ lr\_prob)}\SpecialCharTok{$}\NormalTok{concordance)}
\NormalTok{    cerr\_rsf[i] }\OtherTok{\textless{}{-}} \DecValTok{1} \SpecialCharTok{{-}} \FunctionTok{as.numeric}\NormalTok{(}\FunctionTok{concordance}\NormalTok{(}
      \FunctionTok{Surv}\NormalTok{(test}\SpecialCharTok{$}\NormalTok{time, test}\SpecialCharTok{$}\NormalTok{status) }\SpecialCharTok{\textasciitilde{}}\NormalTok{ rsf\_risk)}\SpecialCharTok{$}\NormalTok{concordance)}
    
    \DocumentationTok{\#\# Step 2g: Clean up}
    \FunctionTok{rm}\NormalTok{(rsf\_mod, rsf\_pred, dat, train, test)}
    \FunctionTok{gc}\NormalTok{()}
\NormalTok{  \}}
  
  \FunctionTok{data.frame}\NormalTok{(}
    \AttributeTok{Scenario     =}\NormalTok{ scenario\_label,}
    \AttributeTok{AUPRC\_LR     =} \FunctionTok{mean}\NormalTok{(auprc\_lr),   }\AttributeTok{AUPRC\_LR\_SD  =} \FunctionTok{sd}\NormalTok{(auprc\_lr),}
    \AttributeTok{AUPRC\_RSF    =} \FunctionTok{mean}\NormalTok{(auprc\_rsf),  }\AttributeTok{AUPRC\_RSF\_SD =} \FunctionTok{sd}\NormalTok{(auprc\_rsf),}
    \AttributeTok{CErr\_LR      =} \FunctionTok{mean}\NormalTok{(cerr\_lr),    }\AttributeTok{CErr\_LR\_SD   =} \FunctionTok{sd}\NormalTok{(cerr\_lr),}
    \AttributeTok{CErr\_RSF     =} \FunctionTok{mean}\NormalTok{(cerr\_rsf),   }\AttributeTok{CErr\_RSF\_SD  =} \FunctionTok{sd}\NormalTok{(cerr\_rsf),}
    \AttributeTok{Events\_mean  =} \FunctionTok{round}\NormalTok{(}\FunctionTok{mean}\NormalTok{(n\_events)),}
    \AttributeTok{Events\_SD    =} \FunctionTok{round}\NormalTok{(}\FunctionTok{sd}\NormalTok{(n\_events))}
\NormalTok{  )}
\NormalTok{\}}

\NormalTok{betas\_A }\OtherTok{\textless{}{-}} \FunctionTok{c}\NormalTok{(}\FloatTok{0.5}\NormalTok{, }\FloatTok{0.4}\NormalTok{, }\FloatTok{0.3}\NormalTok{, }\FloatTok{0.3}\NormalTok{, }\FloatTok{0.2}\NormalTok{, }\FloatTok{0.2}\NormalTok{, }\FloatTok{0.1}\NormalTok{, }\FloatTok{0.1}\NormalTok{, }\FloatTok{0.05}\NormalTok{, }\FloatTok{0.05}\NormalTok{)}
\end{Highlighting}
\end{Shaded}

Note: concordance() uses the survival convention where higher predictor
= longer survival. Since lr\_prob and rsf\_risk are RISK scores (higher
= sooner failure), the raw concordance returns \textless{} 0.5. CErr = 1
- concordance() therefore recovers the true discriminating C-index
(\textgreater{} 0.5 = good).

\textbf{I will need to state this explicitly when reporting these values
in the thesis.}

\subsection{Chunk 3: Run}\label{chunk-3-run}

\begin{Shaded}
\begin{Highlighting}[]
\DocumentationTok{\#\# Step 3: Execute scenarios}
\DocumentationTok{\#\# Delete dgp\_scenario\_results.rds to re{-}run}

\ControlFlowTok{if}\NormalTok{ (}\SpecialCharTok{!}\FunctionTok{file.exists}\NormalTok{(}\StringTok{"dgp\_scenario\_results.rds"}\NormalTok{)) \{}
\NormalTok{  results\_A }\OtherTok{\textless{}{-}} \FunctionTok{run\_scenario}\NormalTok{(}
    \AttributeTok{scenario\_label =} \StringTok{"A: Linear additive, high censoring"}\NormalTok{,}
    \AttributeTok{n              =} \DecValTok{10000}\NormalTok{,}
    \AttributeTok{betas          =}\NormalTok{ betas\_A,}
    \AttributeTok{interactions   =} \ConstantTok{NULL}\NormalTok{,}
    \AttributeTok{censoring\_rate =} \FloatTok{0.99}
\NormalTok{  )}
\NormalTok{  results\_B }\OtherTok{\textless{}{-}} \FunctionTok{run\_scenario}\NormalTok{(}
    \AttributeTok{scenario\_label =} \StringTok{"B: Interactions, moderate censoring"}\NormalTok{,}
    \AttributeTok{n              =} \DecValTok{5000}\NormalTok{,}
    \AttributeTok{betas          =}\NormalTok{ betas\_A,}
    \AttributeTok{interactions   =} \FunctionTok{list}\NormalTok{(}
      \FunctionTok{list}\NormalTok{(}\AttributeTok{i =} \DecValTok{1}\NormalTok{, }\AttributeTok{j =} \DecValTok{3}\NormalTok{, }\AttributeTok{gamma =} \FloatTok{0.5}\NormalTok{),}
      \FunctionTok{list}\NormalTok{(}\AttributeTok{i =} \DecValTok{2}\NormalTok{, }\AttributeTok{j =} \DecValTok{4}\NormalTok{, }\AttributeTok{gamma =} \FloatTok{0.5}\NormalTok{)}
\NormalTok{    ),}
    \AttributeTok{censoring\_rate =} \FloatTok{0.50}
\NormalTok{  )}
\NormalTok{  results\_combined }\OtherTok{\textless{}{-}} \FunctionTok{rbind}\NormalTok{(results\_A, results\_B)}
  \FunctionTok{saveRDS}\NormalTok{(results\_combined, }\StringTok{"dgp\_scenario\_results.rds"}\NormalTok{)}
\NormalTok{\} }\ControlFlowTok{else}\NormalTok{ \{}
\NormalTok{  results\_combined }\OtherTok{\textless{}{-}} \FunctionTok{readRDS}\NormalTok{(}\StringTok{"dgp\_scenario\_results.rds"}\NormalTok{)}
\NormalTok{\}}
\end{Highlighting}
\end{Shaded}

\subsection{Chunk 4: Display}\label{chunk-4-display}

\begin{Shaded}
\begin{Highlighting}[]
\NormalTok{pander}\SpecialCharTok{::}\FunctionTok{pander}\NormalTok{(results\_combined, }\AttributeTok{digits =} \DecValTok{4}\NormalTok{,}
       \AttributeTok{caption =} \StringTok{"Table: DGP Scenario Reproduction — LR vs RSF under controlled conditions"}\NormalTok{)}
\end{Highlighting}
\end{Shaded}

\begin{longtable}[]{@{}
  >{\centering\arraybackslash}p{(\linewidth - 8\tabcolsep) * \real{0.3500}}
  >{\centering\arraybackslash}p{(\linewidth - 8\tabcolsep) * \real{0.1375}}
  >{\centering\arraybackslash}p{(\linewidth - 8\tabcolsep) * \real{0.1750}}
  >{\centering\arraybackslash}p{(\linewidth - 8\tabcolsep) * \real{0.1500}}
  >{\centering\arraybackslash}p{(\linewidth - 8\tabcolsep) * \real{0.1875}}@{}}
\caption{Table: DGP Scenario Reproduction --- LR vs RSF under controlled
conditions (continued below)}\tabularnewline
\toprule\noalign{}
\begin{minipage}[b]{\linewidth}\centering
Scenario
\end{minipage} & \begin{minipage}[b]{\linewidth}\centering
AUPRC\_LR
\end{minipage} & \begin{minipage}[b]{\linewidth}\centering
AUPRC\_LR\_SD
\end{minipage} & \begin{minipage}[b]{\linewidth}\centering
AUPRC\_RSF
\end{minipage} & \begin{minipage}[b]{\linewidth}\centering
AUPRC\_RSF\_SD
\end{minipage} \\
\midrule\noalign{}
\endfirsthead
\toprule\noalign{}
\begin{minipage}[b]{\linewidth}\centering
Scenario
\end{minipage} & \begin{minipage}[b]{\linewidth}\centering
AUPRC\_LR
\end{minipage} & \begin{minipage}[b]{\linewidth}\centering
AUPRC\_LR\_SD
\end{minipage} & \begin{minipage}[b]{\linewidth}\centering
AUPRC\_RSF
\end{minipage} & \begin{minipage}[b]{\linewidth}\centering
AUPRC\_RSF\_SD
\end{minipage} \\
\midrule\noalign{}
\endhead
\bottomrule\noalign{}
\endlastfoot
A: Linear additive, high censoring & 0.03335 & 0.02167 & 0.02151 &
0.007138 \\
B: Interactions, moderate censoring & 0.6298 & 0.01609 & 0.6498 &
0.01904 \\
\end{longtable}

\begin{longtable}[]{@{}
  >{\centering\arraybackslash}p{(\linewidth - 10\tabcolsep) * \real{0.1316}}
  >{\centering\arraybackslash}p{(\linewidth - 10\tabcolsep) * \real{0.1711}}
  >{\centering\arraybackslash}p{(\linewidth - 10\tabcolsep) * \real{0.1447}}
  >{\centering\arraybackslash}p{(\linewidth - 10\tabcolsep) * \real{0.1842}}
  >{\centering\arraybackslash}p{(\linewidth - 10\tabcolsep) * \real{0.1842}}
  >{\centering\arraybackslash}p{(\linewidth - 10\tabcolsep) * \real{0.1842}}@{}}
\toprule\noalign{}
\begin{minipage}[b]{\linewidth}\centering
CErr\_LR
\end{minipage} & \begin{minipage}[b]{\linewidth}\centering
CErr\_LR\_SD
\end{minipage} & \begin{minipage}[b]{\linewidth}\centering
CErr\_RSF
\end{minipage} & \begin{minipage}[b]{\linewidth}\centering
CErr\_RSF\_SD
\end{minipage} & \begin{minipage}[b]{\linewidth}\centering
Events\_mean
\end{minipage} & \begin{minipage}[b]{\linewidth}\centering
Events\_SD
\end{minipage} \\
\midrule\noalign{}
\endhead
\bottomrule\noalign{}
\endlastfoot
0.7152 & 0.05319 & 0.6474 & 0.06304 & 100 & 9 \\
0.7046 & 0.008647 & 0.7257 & 0.008396 & 2074 & 27 \\
\end{longtable}

\end{document}
